\SetRegion{关于本报告}{teal}{compass}
\phantomsection\label{ch1-about}\pdfbookmark[0]{关于本报告}{ch1-bm-about}
\kicker{ABOUT THIS REPORT / 出版与阅读说明}
\PageTitle{关于本报告}{2026 全球开源发展报告}
\lead{以可追溯的证据，理解开放协作的发展条件。}
本报告面向政府与国际组织决策者，同时服务于企业、基金会、研究机构和开源社区。研究关注人工智能时代的参与机会、开放供给、跨主体复用、治理条件与公共价值，既呈现全球格局，也考察不同地区和主体的具体贡献。

\InfoLabel{研究对象与资料基础}
研究对象包括开发者、组织、项目、代码仓库、模型、数据及其协作复用关系。主要使用GitHub与Hugging Face资料，补充平台分别说明覆盖。OpenDigger提供数据、指标与方法；OpenShare是访问和呈现应用。同一上游的不同访问路径不计作独立来源。报告另行固定样本、窗口和版本，具体出处随图列明。

\InfoLabel{统计范围与阅读方法}
主统计期为\textbf{2025年7月1日至2026年6月30日}。历史序列、累计快照和外部研究分别说明时段与对象。开发者与活动、项目与仓库、发布与采用、网络影响与公共成效不互相等同；不同平台或不同样本的数量不未经核对直接合并。

\InfoLabel{目标框架与国际参照}
OSDGs的中文名称为\textbf{开源发展目标}，OSDI为\textbf{开源发展指标体系}。四个目标、十个观察方向与29项候选指标保持不同层级。联合国可持续发展目标是本报告的价值参照，不表示联合国或其他机构为本报告及其指标作出认证或背书。

\InfoLabel{引用、来源与复用}
建议引用：\textbf{《2026全球开源发展报告》编制组．2026全球开源发展报告：第二轮整合评审稿v2.0［R］．2026．}

正文引注可点击进入来源；详细来源记录随配套源码保存。引用时保留对象、时期、单位与来源。第三方内容遵循原许可；指标与数据说明见附录A；模型开放条件见第五章，公共价值与行动建议见第一章、第六章。

\InfoLabel{版本与收录范围}
本稿收录前置内容、六章、附录A与致谢，统一目录、页码与书签，保留竖向全景和项目画像。视觉遵循动态星座v2.0.0。\textbf{第二轮整合评审稿v2.0；2026年11月7日；未公开发布。}
\par\clearpage



